\documentclass[twocolumn]{aastex631}
\begin{document}
\title{A residual reionization ionizing-photon-budget shortfall for star-forming galaxies at z\~{}9 (highC systematic corner)}
\author{NebulaMind Lab (autonomous pipeline)}
\affiliation{NebulaMind Lab, \url{https://lab.nebulamind.net}}
\begin{abstract}
Reionization ionizing-photon-budget reconciliation at z\~{}9: star-forming galaxies require f\_esc=0.662 (+0.620/-0.327) to close the budget (Madau-Dickinson SFRD, log xi\_ion=25.5+/-0.15, clumping C=2-5, JWST-SFRD tail), versus indirect-proxy-inferred f\_esc=0.062 (+0.108/-0.039) from LzLCS O32/beta calibrations. Median delta(required-inferred)=+0.568 dex-frac (16-84\%: +0.229 to +1.187); 97\% of the systematic MC shows a shortfall. Conclusion: a genuine SHORTFALL remains, and the sign holds under both O32 and beta calibrations. The result is bounded by the xi\_ion x clumping x proxy-calibration systematic, not by statistics. Generated autonomously from public data (jwst) via the NebulaMind Lab runner: a bounded, reproducible, descriptive study.
\end{abstract}
\section{Introduction}
Introduction: Recent studies have highlighted a potential crisis in the reionization photon budget, with some suggesting that star-forming galaxies may not produce enough ionizing photons to account for the observed reionization [Muñoz2024]. This issue has been discussed in various contexts, including the role of absorption-dominated reionization and the need for increased demands on ionizing sources [Davies2021], as well as assessments of the galaxy ionizing photon budget at high redshifts [Duncan2015]. To address this problem, we aim to reconcile the reionization ionizing-photon-budget using a literature-anchored approach.
\section{Data and method}
Data and method: Our analysis relies on published values for cosmic star formation rate density (SFRD), ionizing photon production efficiency (xi\_ion), and escape fraction (f\_esc) calibrations. We adopt the Madau \& Dickinson (2014) analytic fitting function for SFRD, while using xi\_ion = 10\^{}25.5 ± 0.15 and f\_esc proxy calibrations from LzLCS O32/beta measurements [Chisholm+22, Flury+22; Simmonds+24]. By combining these values, we calculate the required ionizing photon budget to close the reionization gap at z\~{}9.
\section{Result}
Result: Our calculations show that star-forming galaxies require an escape fraction of f\_esc = 0.662 (+0.620/-0.327) to reconcile the reionization ionizing-photon-budget under the Madau-Dickinson SFRD, log xi\_ion=25.5±0.15, and clumping factor C=2-5. However, indirect-proxy-inferred f\_esc from LzLCS O32/beta calibrations yields a significantly lower value of 0.062 (+0.108/-0.039). This discrepancy results in a median delta of +0.568 dex-frac (16-84\%: +0.229 to +1.187), with 97\% of systematic Monte Carlo simulations indicating a shortfall.
\begin{figure}\centering\includegraphics[width=\columnwidth]{result.png}\caption{A residual reionization ionizing-photon-budget shortfall for star-forming galaxies at z\~{}9 (highC systematic corner)}\end{figure}
\section{Caveats}
Caveats: Our analysis is limited by the reliance on published literature values and calibrations, which may not fully capture the complexities of reionization processes. The use of a single selection method for f\_esc proxy calibration introduces potential biases, as different methods can yield varying results. Additionally, our calculations assume a fixed clumping factor range (C=2-5), which may not accurately represent the true distribution of gas densities in the intergalactic medium. Furthermore, the uncalibrated nature of this measurement means that systematic uncertainties are not fully accounted for, potentially affecting the accuracy of our findings.
\section*{References}
\footnotesize
[Muoz2024] Muñoz et al. (2024). Reionization after JWST: a photon budget crisis?. bibcode 2024MNRAS.535L..37M \par [Davies2021] Davies et al. (2021). The Predicament of Absorption-dominated Reionization: Increased Demands on Ionizing Sources. bibcode 2021ApJ...918L..35D \par [Duncan2015] Duncan et al. (2015). Powering reionization: assessing the galaxy ionizing photon budget at z < 10. bibcode 2015MNRAS.451.2030D \par [Park2022] Park et al. (2022). Calibrating excursion set reionization models to approximately conserve ionizing photons. bibcode 2022MNRAS.517..192P \par [Madau2017] Madau (2017). Cosmic Reionization after Planck and before JWST: An Analytic Approach. bibcode 2017ApJ...851...50M

\end{document}
